%-------------------------
% Cover Letter in LaTeX
% Author: Garv Arora
% JPMorganChase — 2027 Software Engineer Program
%------------------------

\documentclass[letterpaper, 11pt]{article}


\usepackage[hidelinks]{hyperref}
\usepackage[top=0.55in, bottom=0.5in, left=0.85in, right=0.85in]{geometry}
\usepackage[skip=6pt]{parskip}       % space between paragraphs, no indent
\usepackage{setspace}

% Clean, readable line spacing
\setstretch{1.10}

% Remove page numbers
\pagestyle{empty}

%-------------------------------------------
\begin{document}

%----------HEADER----------
{\LARGE \textbf{Garv Arora}}

\vspace{4pt}
{\small Vellore, India $\,\cdot\,$ +91 88008-12254 $\,\cdot\,$ \href{mailto:garvarora0205@gmail.com}{garvarora0205@gmail.com} $\,\cdot\,$ \href{https://linkedin.com/in/gaminization}{linkedin.com/in/gaminization} $\,\cdot\,$ \href{https://github.com/gaminization}{github.com/gaminization}}

\vspace{12pt}
\noindent March 2026

\vspace{4pt}
\noindent Talent Acquisition Team \\
JPMorganChase \\
India

\vspace{12pt}

%----------SALUTATION----------
\noindent Dear Hiring Manager,

\vspace{6pt}

%----------BODY----------

% Para 1: Hook — use their own words, establish identity immediately
\noindent
I'm applying for the 2027 Software Engineer Program at JPMorganChase.
One line in your job description stood out to me: \textit{``technologies like Robotics.''}
That's not a throwaway qualifier for me --- it's what I've spent the last two years
actually building, testing under competition conditions, and iterating on when it breaks.

% Para 2: Technical proof + leadership
\vspace{6pt}
\noindent
At Team Vyadh (SEDS VIT), I architected the autonomous navigation stack for our Mars
rover: ROS2, RTAB-Map Visual SLAM, Intel RealSense D455, and an Extended Kalman Filter
sensor fusion system enabling sub-10\,cm autonomous positioning.
Our team placed \textbf{12th globally at IRC 2025} and \textbf{16th at IRC 2026}
among 100+ international teams.
More recently I was promoted to R\&D Lead at SEDS VIT --- the apex body of SEDS India
--- overseeing 50+ engineers across two national competition teams.
Working within large, technically diverse teams toward a hard deadline is something
I've done under real pressure, not just described in a course project.

% Para 3: Force for Good alignment — this is the differentiator
\vspace{6pt}
\noindent
I also built an mmWave-based earthquake survivor detection robot: an LD2410C radar
interfaced with ROS2, capable of detecting breathing signatures through rubble and
plotting survivor coordinates on a live occupancy map.
When I read about JPMorganChase's Force for Good program, this project came to mind
immediately --- using engineering to solve problems that matter for people who need it most
is exactly the kind of work I want to keep doing.

% Para 4: SDLC / engineering practice — maps directly to JD requirements
\vspace{6pt}
\noindent
On the software engineering side, my internship at Wissen Baum Engineering Solutions
gave me direct hands-on experience with the practices your JD calls out:
I built a BDD automation framework in Python, parallelised GitLab CI/CD pipelines
integrating Playwright and Cypress, and replaced manual QA with a pandas/NumPy
validation script processing 10,000+ rows per run.
I understand agile delivery and code quality enforcement --- not as concepts,
but as things I've shipped and maintained.

% Para 5: Closing — concise, confident, no fluff
\vspace{6pt}
\noindent
I'm graduating in July 2027 with a B.Tech in Computer Science and Engineering
(IoT specialisation) from VIT Vellore.
I'm looking for a program where I can contribute technically from day one and grow
alongside engineers who take the craft seriously.
The SEP's structure --- foundational development, career mobility, global exposure ---
is the environment I want to start my career in, and I'm ready to earn my place in it.

\vspace{6pt}
\noindent
I'd welcome the chance to speak further.

%----------SIGN-OFF----------
\vspace{14pt}
\noindent Warm regards,

\vspace{10pt}
\noindent \textbf{Garv Arora}

\end{document}
